% ===== PRÉAMBULE CANONIQUE — à coller VERBATIM en tête de CHAQUE .tex =====
\documentclass[11pt,a4paper]{article}
\usepackage[utf8]{inputenc}\usepackage[T1]{fontenc}\usepackage{lmodern}
\usepackage[french]{babel}
\usepackage{amsmath,amssymb,mathtools}\usepackage{icomma}
\usepackage{siunitx}\sisetup{locale=FR}  % \qty{}{}, \si{}, \ang{}
\usepackage{xcolor}\usepackage{colortbl}
\usepackage{tikz}\usepackage{pgfplots}\pgfplotsset{compat=1.18}
\usepackage[siunitx,europeanresistors]{circuitikz} % circuits (élec.)
\usepackage[most]{tcolorbox}\usepackage{enumitem}\usepackage{array,booktabs}
\usepackage[a4paper,margin=2.2cm]{geometry}
\usetikzlibrary{arrows.meta,calc,angles,quotes,positioning,patterns,%
  backgrounds,shapes.geometric,decorations.pathmorphing,decorations.markings,intersections}
\definecolor{primary}{RGB}{222,49,99}\definecolor{primarydark}{RGB}{150,22,55}
\definecolor{defcol}{RGB}{0,114,178}\definecolor{methcol}{RGB}{52,92,148}
\definecolor{excol}{RGB}{0,135,90}\definecolor{attncol}{RGB}{200,40,55}
\tcbset{boxbase/.style={enhanced,breakable,boxrule=0.9pt,arc=2.5pt,
  left=3mm,right=3mm,top=2.3mm,bottom=2.3mm,fonttitle=\bfseries,coltitle=white}}
% --- boîtes TUTORIEL (noms EXACTS, ne pas renommer) ---
\newtcolorbox{cle}[1][]{boxbase,colback=defcol!6,colframe=defcol,
  title={\textsc{À retenir}\ifx\relax#1\relax\relax\else\textmd{ : \itshape #1}\fi}}
\newtcolorbox{methode}[1][]{boxbase,colback=methcol!7,colframe=methcol,sharp corners,
  title={\textsc{Méthode}\ifx\relax#1\relax\relax\else\textmd{ --- \itshape #1}\fi}}
\newtcolorbox{exemplebox}[1][]{boxbase,colback=excol!5,colframe=excol,
  title={\textsc{Exemple}\ifx\relax#1\relax\relax\else\textmd{ : \itshape #1}\fi}}
\newtcolorbox{piege}[1][]{boxbase,colback=attncol!6,colframe=attncol,
  title={\textsc{Piège}\ifx\relax#1\relax\relax\else\textmd{ : \itshape #1}\fi}}
\newtcolorbox{remarque}[1][]{boxbase,colback=black!4,colframe=black!45,
  title={\textsc{Remarque}\ifx\relax#1\relax\relax\else\textmd{ : \itshape #1}\fi}}
% --- boîtes EXERCICE / CORRIGÉ (noms EXACTS, auto-comptées) ---
\newtcolorbox[auto counter]{exo}[1][]{boxbase,colback=primary!4,colframe=primary,
  title={\textsc{Exercice}~\thetcbcounter\ifx\relax#1\relax\relax\else\textmd{ --- \itshape #1}\fi}}
\newtcolorbox[auto counter]{corrige}[1][]{boxbase,colback=excol!4,colframe=excol,
  title={\textsc{Corrigé}~\thetcbcounter\ifx\relax#1\relax\relax\else\textmd{ --- \itshape #1}\fi}}
\newcommand{\vv}[1]{\vec{#1}}\newcommand{\ds}{\displaystyle}
\newcommand{\R}{\mathbb{R}}\newcommand{\N}{\mathbb{N}}\newcommand{\Z}{\mathbb{Z}}
% (un \usepackage{fancyhdr}+en-tête est possible ; il est ignoré côté web)
% =========================================================================
\begin{document}
\sisetup{per-mode=symbol}

\begin{exo}[énergie dissipée par les frottements de l'air]
Une pierre de masse $m=\qty{5}{\kilogram}$ est lâchée sans vitesse initiale et tombe d'une hauteur
$h=\qty{10}{\metre}$. À cause de la résistance de l'air, elle n'arrive au sol qu'à
$v=\qty{12}{\metre\per\second}$ ($g=\qty{10}{\metre\per\second\squared}$).
\begin{enumerate}[label=\alph*)]
  \item Calcule l'énergie potentielle de pesanteur perdue ($mgh$) et l'énergie cinétique réellement
        acquise ($\tfrac12 m v^2$).
  \item Déduis-en l'énergie dissipée par les frottements de l'air.
\end{enumerate}
\end{exo}

\begin{exo}[une descente avec frottement]
Un skieur de masse $m=\qty{60}{\kilogram}$ descend une piste et perd une dénivellation
$\Delta h=\qty{20}{\metre}$. À cause des frottements, il arrive en bas à $v=\qty{15}{\metre\per\second}$
($g=\qty{10}{\metre\per\second\squared}$).
\begin{enumerate}[label=\alph*)]
  \item Quelle vitesse aurait-il eue \emph{sans} frottement ? (utilise $v=\sqrt{2g\,\Delta h}$)
  \item Calcule l'énergie mécanique dissipée par les frottements ($mg\,\Delta h-\tfrac12 m v^2$).
\end{enumerate}
\end{exo}

\begin{exo}[vrai ou faux : conversions d'énergie d'une balle lancée]
Une balle de masse $m=\qty{0,2}{\kilogram}$ est lancée du sol verticalement vers le haut à
$v_0=\qty{10}{\metre\per\second}$ (sans frottement, $g=\qty{10}{\metre\per\second\squared}$). Indique si
chaque affirmation est \textbf{vraie} ou \textbf{fausse}, en justifiant.
\begin{enumerate}[label=\alph*)]
  \item L'énergie mécanique de la balle vaut \qty{10}{\joule} pendant tout le mouvement.
  \item La hauteur maximale atteinte est \qty{10}{\metre}.
  \item À une hauteur de \qty{2,5}{\metre}, l'énergie potentielle vaut \qty{5}{\joule} et l'énergie
        cinétique \qty{5}{\joule}.
  \item Lorsque la vitesse s'annule, l'énergie potentielle a augmenté de \qty{10}{\joule} par rapport
        au sol.
\end{enumerate}
\end{exo}

\begin{exo}[vrai ou faux : l'indépendance de la masse]
On néglige tous les frottements. Indique si chaque affirmation est \textbf{vraie} ou \textbf{fausse},
en justifiant.
\begin{enumerate}[label=\alph*)]
  \item Deux objets de masses différentes lâchés de la même hauteur arrivent en bas à la même vitesse.
  \item Deux objets lancés verticalement avec la même vitesse initiale montent à la même hauteur,
        quelle que soit leur masse.
  \item Sur une même piste, le skieur le plus lourd arrive en bas avec une vitesse plus grande que le
        plus léger.
  \item La hauteur maximale d'un lancer vertical dépend de la masse de l'objet.
\end{enumerate}
\end{exo}

\begin{exo}[capstone : de la pente au ressort]
Un bloc de masse $m=\qty{2}{\kilogram}$ part du repos en haut d'une piste \emph{sans frottement}, à une
hauteur $h=\qty{1,6}{\metre}$. En bas, la piste devient horizontale et le bloc vient comprimer un
ressort de raideur $k=\qty{1600}{\newton\per\metre}$. On prend $g=\qty{10}{\metre\per\second\squared}$.
\begin{center}
\begin{tikzpicture}[>=Stealth,scale=0.95]
  \draw[primary,line width=1.3pt] (0,2.4) .. controls (1,2.4) and (1.4,0) .. (2.6,0) -- (6.2,0);
  \draw[black!70,line width=1.2pt] (6.2,0)--(6.2,1.4);
  \draw[defcol,line width=1pt] (5.0,0.5) decorate[decoration={coil,aspect=0.6,segment length=5pt,amplitude=4pt}]{-- (6.1,0.5)};
  \fill[excol] (0.1,2.4) circle (0.13); \node[above] at (0.3,2.5){\small $m$};
  \draw[<->,black!55,dashed] (-0.4,0)--(-0.4,2.4) node[midway,left]{\small $h=\qty{1,6}{\metre}$};
  \draw[dashed,black!35] (-0.4,0)--(2.6,0);
  \node[black!70] at (4.2,0.9){\small $k$};
\end{tikzpicture}
\end{center}
\begin{enumerate}[label=\alph*)]
  \item Calcule l'énergie potentielle de pesanteur du bloc en haut (référence en bas). Que devient-elle
        en bas de la piste ?
  \item À la compression maximale, toute l'énergie est stockée dans le ressort ($mgh=\tfrac12 k x^2$).
        Calcule la compression maximale $x$.
\end{enumerate}
\end{exo}
\end{document}
